% ===================================================================
%  تابلوی شمارهٔ ۹ — کوه و شهر آب‌گرم. --- جنگل و شکار.
%
%  Traduction, faite en 2026, en persan d'aujourd'hui, sur l'ido de
%  texto/io. Regles de la colonne : voir 10-tablo-01.tex. Deux
%  scenes, deux numerotations. Un bloc marque « ¶2 » par ordo.py.
%
%  LA MONTAGNE EST LE TABLEAU LE PLUS IRANIEN DU LIVRET, ET LES DEUX
%  SEULES COLONNES QUI PUISSENT LE DIRE SANS EMPRUNTER SONT CELLE-CI
%  ET LA JAVANAISE. کوه la montagne, قله le sommet — non, celui-ci est
%  arabe —, تپه la colline, دره la vallee (25), درّه et تنگه le
%  defile (27), آبشار la cascade (10), چشمه la source (20), رود la
%  riviere, یخچال le glacier (5), بهمن l'avalanche, غار la grotte
%  (56), صخره la roche (45), سنگ la pierre, معدن la mine (28), کان le
%  gisement, کوره le haut-fourneau (29), آهن le fer.
%
%  ET بهمن EST LE MOT LE PLUS SURPRENANT DU RELEVE. Il nomme
%  l'avalanche en persan, et c'est aussi le nom du onzieme mois du
%  calendrier iranien — celui de fin janvier, quand les avalanches
%  tombent sur l'Alborz. Le javanais du tableau 9 avait du composer
%  « longsoran salju », glissement de neige, faute de mot ; le persan
%  en a un si ordinaire qu'il a donne son nom a un mois. La meme
%  planche, le meme objet, et une langue qui compte ses saisons par
%  lui.
%
%  ET LE VOLCAN (6) TRANCHE DANS L'AUTRE SENS, POUR UNE FOIS. Java
%  est une ile de volcans et le javanais avait kawah le cratere,
%  lahar la coulee, wedhus gembel la nuee — un vocabulaire complet et
%  a lui. L'Iran a le Damavand, mais il dort depuis sept mille ans :
%  le persan dit آتشفشان, « le feu-crachant », un compose transparent
%  du XXe siecle, et آتشفشان خاموش pour le volcan eteint de la
%  planche. Une langue nomme richement ce qu'elle craint tous les
%  ans, et decrit ce qu'elle n'a vu que dans un livre. Le partage est
%  exactement l'inverse de celui de la neige au tableau 7.
%
%  UNE FAUTE D'ORDRE, LA CINQUIEME DE LA COLONNE, ET C'EST LA
%  PREMIERE OU L'IDO ETAIT LE PLUS NATUREL. Au bloc t09-c2-09-1 le
%  cerf (50) avait ete pose avant le taillis (54) dont il sort, alors
%  que l'ido dit « vidis debushar ek kopso belega cervulo » — la
%  source elle-meme pose la source du mouvement avant le mobile. Le
%  persan le fait aussi bien : « از بیشهٔ کوتاه گوزنی بیرون می‌زند ».
%  Cinq fautes en neuf tableaux, toutes de placement, et pas une de
%  grammaire.
%
%  LA FORET, ELLE, EST LE SEUL ENDROIT DE CETTE COLONNE OU LE PERSAN
%  MANQUE DE MOTS, ET LA RAISON EST UN CLIMAT. Le plateau iranien est
%  sec : les grandes forets de hetres et de chenes de la seconde scene
%  n'existent que sur la rive caspienne. راش le hetre (23), بلوط le
%  chene (29), صنوبر le peuplier (48), نارون l'orme (47), کاج le pin,
%  سرو le cypres, نمدار le tilleul — les sept sont persans, mais ce
%  sont les arbres du Hyrcanien, et le tableau les nomme tous parce
%  que cette foret-la existe. Le persan n'a manque de rien : c'est le
%  releve qui s'attendait a ce qu'il manque.
%
%  ET LES BETES DE LA FORET SONT PERSANES SANS EXCEPTION : خرس l'ours
%  (44), گراز le sanglier (21), گوزن le cerf (50), آهو le chevreuil
%  (20), روباه le renard, خرگوش le lievre (18), سنجاب l'ecureuil
%  (61), مار le serpent (7), عنکبوت — non, celui-ci est arabe —,
%  مورچه la fourmi, زنبور, پروانه, قورباغه, لاک‌پشت. Le seul emprunt
%  du bestiaire est l'araignee, et c'est l'exception qui rappelle la
%  regle du tableau 2 : ce qu'on nomme dans un livre est arabe, ce
%  qu'on rencontre est persan.
% ===================================================================

%%K t09-apar-1 apar
\VUsaut{4.0mm}
\VUtitre{15.0pt}{60}{تابلوی شمارهٔ ۹}
\VUsaut{3.0mm}

%%K t09-tit-1 sub
\VUsaut{3.0mm}
\VUcentre{12.0pt}{0}{\textit{صحنهٔ اول.}}
\VUsaut{3.0mm}

%%K t09-tit-2 sub
\VUsaut{3.0mm}
\VUpk{11.4pt}{40}{کوه. --- شهر آب‌گرم.}
\VUsaut{3.0mm}

%%K t09-c1-01-1 p
1. --- هشت روز است که در پراتس هستم. نمی‌توانم همهٔ شگفتی خود را بیان
کنم که در برابر \VUgras{رشته‌کوه} \textsuperscript{(1)} شگرف پیرنه
احساس کردم، که \VUgras{ستیغ} \textsuperscript{(2)} آن بر آسمان آبی
مانند حاشیه‌ای غول‌آسا می‌نماید.

%%K t09-c1-02-1 p
2. --- گمان نمی‌کردم که \VUgras{قله‌های} \textsuperscript{(3)} پیرنه به
\VUgras{بلندی} چنین بزرگی برسند، و در برابر \VUgras{یخچال‌های}
\textsuperscript{(5)} بسیار زیبا و \VUgras{نوک‌های کوه}
\textsuperscript{(4)} برفی که \VUgras{بهمن‌های} چنین هراس‌انگیز بر
آنها فرو می‌غلتند، سخت شگفت‌زده ماندم.

%%K t09-tit-3 sub
\VUsaut{3.0mm}
\VUpk{10.2pt}{40}{چشم‌انداز کوهستان.}
\VUsaut{3.0mm}

%%K t09-c1-03-1 p
3. --- همان فردای رسیدنم، به \VUgras{آتشفشان} \textsuperscript{(6)}
خاموش کهنی رفتم که نزدیک پراتس است. بر کناره‌های خشک و برهنهٔ
\VUgras{دهانه}، هنوز نشانه‌های گذر \VUgras{گدازه} را خوب می‌بینیم که
چند سده پیش بر \VUgras{دامنهٔ کوه} \textsuperscript{(7)} روان شد.
توانستم بر یکی از \VUgras{شیب‌ها} چند تکه از آن گدازه بیابم.

%%K t09-c1-04-1 p
4. --- پراتس بر مرز جای گرفته است، و \VUgras{دژ} \textsuperscript{(8)}
که از \VUgras{تنگهٔ} \textsuperscript{(27)} جداکنندهٔ فرانسه از
اسپانیا نگهبانی می‌کند، سراسر کاخ‌های کهن کرانهٔ راین را به یاد
می‌آورد که گویی از فراز صخرهٔ خود در کمین \VUgras{شکار} نشسته‌اند.

%%K t09-c1-05-1 p
5. --- \VUgras{عقاب} \textsuperscript{(9)} بسیار بزرگی که آشیانه‌اش نه
چندان دور از \VUgras{قلعه} \textsuperscript{(8)} است، بارها توجه مرا
جلب می‌کند.

آن \VUgras{پرندهٔ شکاری} که \VUgras{منقارش} نیرومند است، بارها برای
جوجه‌هایش بره یا بزغاله‌ای می‌آورد که با \VUgras{چنگال‌های} خود نگه
داشته است. \VUgras{بال‌های} او چندان بزرگ‌اند که می‌تواند با بار
سنگین خود دیرگاهی سُر بخورد.

%%K t09-tit-4 sub
\VUsaut{3.0mm}
\VUpk{10.2pt}{40}{گردشگر.}
\VUsaut{3.0mm}

%%K t09-c1-06-1 p
6. --- در آنجا خوشایندترین گردش، رفتن به آبشار \textsuperscript{(10)}
« کاو » است. \VUgras{آبشار} گردابی فرو می‌ریزد و سپس مانند
\VUgras{جویبار} زیبایی در \VUgras{جنگل کاج} \textsuperscript{(11)}
ناپدید می‌شود که در باختر شهر است. از این جنگل تا \VUgras{ایستگاه
هواشناسی} \textsuperscript{(12)} بالا رفتیم، و از فراز
\VUgras{دیدگاه}، همهٔ \VUgras{چشم‌انداز} منطقه را دیدیم.
\VUgras{منظره} شگفت‌انگیز است، و \VUgras{طبیعت} گویی همهٔ گنج‌هایی را
که دارد به این سرزمین می‌بخشد.

%%K t09-c1-07-1 p
7. --- برای پایین آمدن، از \VUgras{تله‌کابین} \textsuperscript{(13)}
بهره بردیم و در \VUgras{ایستگاه} کوچکی نزدیک \VUgras{ویرانه‌های}
\textsuperscript{(14)} \VUgras{کاخ فئودالی} کهنی ایستادیم که روزگاری
سروران پراتس در آن می‌نشستند.

%%K t09-c1-08-1 p
8. --- کاخ بینوا! گمان می‌رود چه می‌اندیشد هنگامی که در پایین
\VUgras{شهر آب‌گرم} \textsuperscript{(15)} خوش‌آب‌ورنگ را می‌نگرد که
اکنون \VUgras{ایستگاه آب‌گرم} مد روز است؟

%%K t09-c1-08-2 p
\VUgras{آب‌وهوا} چنان خوشایند و \VUgras{آب معدنی} چنان کارگر است که
\VUgras{درمانی} که در \VUgras{آسایشگاه} \textsuperscript{(17)}
به کار می‌رود لذت راستین است.

%%K t09-tit-5 sub
\VUsaut{3.0mm}
\VUpk{10.2pt}{40}{شهر آب‌گرم دلپذیر.}
\VUsaut{3.0mm}

%%K t09-c1-09-1 p
9. --- وانگهی، برای خوشایند کردن اقامت همه کار کرده‌اند.
\VUgras{پل بلند} \textsuperscript{(16)} بزرگی ساخته شده است تا راه‌آهن
بگذرد ؛ \VUgras{پارک} \textsuperscript{(18)} \VUgras{مجموعهٔ آب‌گرم}
که در آن \VUgras{کنسرت می‌دهند}، به شیوهٔ دلربایی آراسته شده است. چند
« \VUgras{ویلا} » \textsuperscript{(19)} خوش‌دوخت گرد کازینو
\textsuperscript{(21)} ساخته شده‌اند، و \VUgras{جادهٔ}
\textsuperscript{(22)} بسیار خوب نگه‌داشته‌شده رفت‌وآمد را بسیار آسان
می‌کند.

%%K t09-c1-10-1 p
10. --- \VUgras{خاکریز} \textsuperscript{(23)} ساده‌ای \VUgras{راه}
\textsuperscript{(20)} را از \VUgras{درهٔ} \textsuperscript{(25)}
اصلی جدا می‌کند که در ژرفای آن \VUgras{دریاچهٔ} \textsuperscript{(26)}
کوچک زیبایی خفته است و \VUgras{جویبار} \textsuperscript{(24)} زلالی را
سیراب می‌کند.

%%K t09-c1-11-1 p
11. --- در سوی دیگر دره، \VUgras{معدن آهن} \textsuperscript{(28)}
بهره‌برداری می‌شود. بارها رفتم تا \VUgras{معدنچیان} را ببینم که به
\VUgras{چاه} فرو می‌روند، و حتی به \VUgras{دالانی} درآمدم که روز خود
را در آن می‌گذرانند و \VUgras{سنگ معدن} را بیرون می‌کشند که
\VUgras{فلز} در آن است.

%%K t09-c1-12-1 p
12. --- \VUgras{کورهٔ ذوب} بزرگی نزدیک معدن ساخته شده است تا گنج‌هایی
را که از آن بیرون می‌کشند بی‌درنگ به کار برند. \VUgras{کورهٔ بلند}
\textsuperscript{(29)} که با \VUgras{زغال سنگ} فراوان اینجا گرم
می‌شود، به دست آوردن \VUgras{چدن} را زود و ارزان شدنی می‌کند.

%%K t09-c1-13-1 p
13. --- نه چندان دور از معدن، \VUgras{معدن سنگ} \textsuperscript{(30)}
می‌کنند. \VUgras{سنگ‌کار} پیر نه \VUgras{گرانیت}، نه \VUgras{سنگ سماق}
و نه حتی \VUgras{ماسه‌سنگ} را با \VUgras{کلنگ} خود می‌برد، بلکه تنها
\VUgras{سنگ تراشیده} برای ساختن خانه.

%%K t09-c1-14-1 p
14. --- یکی از زیباترین آرایه‌های آن سرزمین \VUgras{کاج‌ها}
\textsuperscript{(31)} هستند. همه‌جا در ژرفای \VUgras{تنگه}
\textsuperscript{(32)}، بر کنارهٔ \VUgras{سیلاب} \textsuperscript{(33)}،
\VUgras{پرتگاه} \textsuperscript{(34)} یا شیب تند، سوزن‌های باریکشان
سبزی می‌بخشد.

%%K t09-tit-6 sub
\VUsaut{3.0mm}
\VUpk{10.2pt}{40}{پیاده‌روی.}
\VUsaut{3.0mm}

%%K t09-c1-15-1 p
15. --- از تعطیلاتم بهره می‌برم تا \VUgras{دور} گردش کنم.
\VUgras{راه‌هایی} \textsuperscript{(35)} که در کوه پیچ‌وخم می‌خورند،
این را شدنی می‌کنند که بی‌آنکه دوری را احساس کنی چند \VUgras{کیلومتر}
و بارها چند \VUgras{ده کیلومتر} را بپیمایی.

%%K t09-c1-15-2 p
\VUgras{سنگ‌های مرزی} \textsuperscript{(36)} و \VUgras{تیرک‌های
راهنما} \textsuperscript{(37)} راه‌ها را نشانه گذاشته‌اند، جایی که
بارها با \VUgras{کوهنشین} \textsuperscript{(38)} جوانی روبه‌رو
می‌شویم که \VUgras{خورجین} \textsuperscript{(40)} بر پهلو دارد و
\VUgras{موش کوهی} \textsuperscript{(39)} می‌برد، یا با چند
\VUgras{کولی} \textsuperscript{(41)} که \VUgras{پوستین}
\textsuperscript{(42)} آنان با \VUgras{کلاه} \textsuperscript{(43)}
کهنه و پاره‌شان تضادی شگفت دارد. او \VUgras{خرسی}
\textsuperscript{(44)} تربیت‌شده را با \VUgras{زنجیر} می‌برد.

%%K t09-tit-7 sub
\VUsaut{3.0mm}
\VUpk{10.2pt}{40}{کوهنوردی.}
\VUsaut{3.0mm}

%%K t09-c1-16-1 p
16. --- تقریباً هر روز با \VUgras{راهنما} \textsuperscript{(48)} برای
\VUgras{صعود} می‌روم، و بسیار سربلندم هنگامی که با \VUgras{کوله‌پشتی}
\textsuperscript{(47)} بر پشت و « \VUgras{عصای کوهنوردی} » در دست،
مانند \VUgras{گردشگر} \textsuperscript{(46)} راستین به \VUgras{صخره‌ها}
\textsuperscript{(45)} یورش می‌برم.

%%K t09-c1-17-1 p
17. --- از فراز کوه، ساکنان دشت بزرگ‌تر از مورچه نمی‌نمایند ؛
\VUgras{گلهٔ گوسفند} \textsuperscript{(50)} که در \VUgras{چراگاه}
\textsuperscript{(49)} می‌چرد، اسباب‌بازی کودکان می‌نماید. \VUgras{کاج}
\textsuperscript{(51)} یا حتی \VUgras{خانهٔ چوبی} \textsuperscript{(52)}
به سختی چون لکه‌ای بر سبزه می‌نماید.

%%K t09-c1-18-1 p
18. --- در این گردش‌ها، بارها بر \VUgras{راه باریک}
\textsuperscript{(53)} با \VUgras{شکارچی بز کوهی} \textsuperscript{(54)}
روبه‌رو می‌شویم که جانوری را که همین حالا کشته است بر دوش می‌برد.

%%K t09-c1-19-1 p
19. --- در گیاه‌نامهٔ خود شاخهٔ بسیار چشمگیری از \VUgras{سرخس}
\textsuperscript{(55)} و شاخهٔ کوچک صورتی زیبایی از \VUgras{خلنگ}
\textsuperscript{(57)} گذاشتم که در ورودی \VUgras{غار}
\textsuperscript{(56)} کوچکی در واپسین گردشم چیدم.

%%K t09-tit-8 sub
\VUsaut{3.0mm}
\VUcentre{12.0pt}{0}{\textit{صحنهٔ دوم.}}
\VUsaut{3.0mm}

%%K t09-tit-9 sub
\VUsaut{3.0mm}
\VUpk{11.4pt}{40}{جنگل. --- شکار.}
\VUsaut{3.0mm}

%%K t09-c2-01-1 p
1. --- دیروز روز بسیار دلپذیری را در جنگل گذراندم. کوششی که میان
جانوران و \VUgras{حشره‌های} \textsuperscript{(5)} ساکن آن فرمان
می‌راند، مرا بسیار به شوق آورد.

%%K t09-c2-01-2 p
\VUgras{عنکبوتی} \textsuperscript{(1)} که \VUgras{تار}
\textsuperscript{(2)} خود را میان دو شاخه می‌تند تا \VUgras{مگسی}
\textsuperscript{(3)} را که می‌گذرد بگیرد، \VUgras{حلزونی}
\textsuperscript{(4)} که با رنج صدف خود را می‌برد، سوسک شاخداری که
شکار خود را می‌جوید، مورچهٔ کوشا که نزدیک \VUgras{لانهٔ مورچه}
\textsuperscript{(6)} بسیار پرشور است — همهٔ این کوچک‌ها می‌رفتند،
می‌آمدند، با کوششی خارق‌العاده کار می‌کردند.

%%K t09-c2-02-1 p
2. --- \VUgras{ماری} \textsuperscript{(7)} که در نگاه نخست پنداشتم
\VUgras{افعی} است، اما جز \VUgras{مار آبی} ساده‌ای نبود، زیر تنهٔ
درختی پنهان شده بود و گاه‌گاه سر باریک خود را بیرون می‌آورد که همیشه
گویی آمادهٔ گزیدن بود. جلوِ من \VUgras{راب} \textsuperscript{(8)}
بزرگی سنگین می‌خزید، پس از آنکه یکی از \VUgras{توت‌فرنگی‌های}
\textsuperscript{(11)} خوشمزه را خورده بود که فراوان آنجا می‌رویند.

%%K t09-c2-03-1 p
3. --- زمین با \VUgras{گل‌های جنگلی} فرش شده بود ؛ \VUgras{پامچال‌ها}
\textsuperscript{(9)}، \VUgras{پروانش‌ها} \textsuperscript{(10)}، و
گیاهان بسیار دیگری که نمی‌شناختم، میان \VUgras{دسته‌های علف}
\textsuperscript{(12)} شکفته بودند.

%%K t09-c2-04-1 p
4. --- در آن منطقه، \VUgras{دنبلان} تقریباً به همان اندازهٔ
\VUgras{قارچ} \textsuperscript{(14)} فراوان است، و \VUgras{بچه‌خوکی}
\textsuperscript{(13)} از آنِ جنگلبان، شکم‌پرستانه می‌جست که آیا
خواهد توانست چند تایی از آنها را بیابد.

%%K t09-tit-10 sub
\VUsaut{3.0mm}
\VUpk{10.2pt}{40}{شکار قاچاق.}
\VUsaut{3.0mm}

%%K t09-c2-05-1 p
5. --- در \VUgras{پایان} صحنه‌ای بودم که مرا بسیار سرگرم کرد. به یاری
بویایی \VUgras{سگ باست} \textsuperscript{(15)} خود، \VUgras{جنگلبان}
\textsuperscript{(16)} همین حالا \VUgras{شکارچی قاچاق}
\textsuperscript{(22)} را غافلگیر کرده بود، درست هنگامی که او آماده
می‌شد \VUgras{شکاری} \textsuperscript{(17)} را که کشته بود با خود ببرد.
شکارچی قاچاق که رها کردن شکار را بر دستگیر شدن ترجیح داده بود، گریخته
بود، و \VUgras{خرگوش} \textsuperscript{(18)}، \VUgras{ماده‌گوزن}
\textsuperscript{(19)}، \VUgras{آهو} \textsuperscript{(20)} و
\VUgras{گراز} \textsuperscript{(21)} روی علف افتاده بودند و جنگلبان را
شاد می‌کردند.

%%K t09-c2-06-1 p
6. --- گوناگونی درختانی که جنگل در بر دارد شگفت‌انگیز است.
\VUgras{راش‌ها} \textsuperscript{(23)} و \VUgras{توس‌ها}
\textsuperscript{(26)}، \VUgras{کاج‌ها} که \VUgras{صمغ}
\textsuperscript{(27)} می‌دهند، \VUgras{بلوط‌ها} \textsuperscript{(29)}
و بسیار دیگر، شاخ‌وبرگ خود را در هم می‌آمیزند.

%%K t09-c2-06-2 p
\VUgras{پیچک} \textsuperscript{(24)} که به تنهٔ درختان بلند می‌چسبد،
\VUgras{بیشه} \textsuperscript{(25)} را با سبزی درخشانی رنگ می‌زند که
به خوشی با رنگ روشن‌تر و سبز کم‌رنگ \VUgras{خزه} \textsuperscript{(31)}
می‌آمیزد، جایی که \VUgras{دارکوب سبز} \textsuperscript{(30)} حشره
می‌جوید.

%%K t09-tit-11 sub
\VUsaut{3.0mm}
\VUpk{10.2pt}{40}{چشم‌انداز جنگل.}
\VUsaut{3.0mm}

%%K t09-c2-07-1 p
7. --- بلوط‌های کهن مرا شیفته کردند. \VUgras{ریشه‌های}
\textsuperscript{(32)} کلفت و پیچیدهٔ آنها که نیمی از زمین بیرون
آمده‌اند، نمایی شکوهمند از نیرو و توان به آنها می‌دهند، و از خود
می‌پرسم چگونه \VUgras{مالک} \textsuperscript{(35)} می‌تواند به بریدن
آنها و سپردنشان به \VUgras{ارهٔ} \textsuperscript{(34)}
\VUgras{هیزم‌شکن} \textsuperscript{(33)} تن دهد. \VUgras{تنه‌های}
\textsuperscript{(36)} بینوا که از \VUgras{پوست} \textsuperscript{(37)}
خود برهنه شده‌اند یا با \VUgras{تبر} \textsuperscript{(38)}
\VUgras{کارگر روزمزد} \textsuperscript{(39)} شکافته شده‌اند، مرا
اندوهگین می‌کنند.

%%K t09-c2-08-1 p
8. --- نزدیک آنجا، \VUgras{زغال‌ساز} \textsuperscript{(41)} پیری با
\VUgras{بیل} خود \VUgras{تودهٔ} \textsuperscript{(42)} زغالی آماده
کرده است، و پیرزنی \VUgras{بستهٔ} \textsuperscript{(40)} \VUgras{هیزم
خشک} را از شاخه‌های درختان بریده‌شده با خود برده است.

%%K t09-tit-12 sub
\VUsaut{3.0mm}
\VUpk{10.2pt}{40}{شکار.}
\VUsaut{3.0mm}

%%K t09-c2-09-1 p
9. --- می‌خواستم چند لحظه در \VUgras{خانهٔ جنگلبان}
\textsuperscript{(46)} بیاسایم، اما هنگامی که نزدیک \VUgras{برکهٔ}
\textsuperscript{(43)} کوچکی رسیدم که \VUgras{قوی} \textsuperscript{(45)}
زیبایی روی آن شنا می‌کرد، دیدم که \VUgras{سیلاب} \textsuperscript{(44)}
تازه‌ای راه را به \VUgras{باتلاق} راستینی دگرگون کرده است. ناچار شدم
راه را کج کنم، و در گذر از نزدیک \VUgras{درخت سدر}
\textsuperscript{(49)}، \VUgras{صنوبرها} \textsuperscript{(48)} و
\VUgras{نارونی} \textsuperscript{(47)} که بر کنارهٔ جنگل‌اند، دیدم که
از \VUgras{بیشهٔ کوتاه} \textsuperscript{(54)} \VUgras{گوزن}
\textsuperscript{(50)} بسیار زیبایی بیرون می‌زند، دنبال‌شده از \VUgras{دستهٔ
سگ شکاری} \textsuperscript{(51)} پیگیری که \VUgras{بوق}
\textsuperscript{(53)} \VUgras{شکارچیان سوار} \textsuperscript{(52)} نیز
آن را برانگیخته بود. جانور بینوا سراسر از توان افتاده بود. آمد و چند
گام دورتر از من مرده افتاد.

%%K t09-c2-10-1 p
10. --- پسر جنگلبان که هیاهوی \VUgras{شکار با سگ} او را کشیده بود، از
خانه بیرون آمد و ما دو تن دوباره به جنگل بازگشتیم. او \VUgras{زاغی}
\textsuperscript{(56)} را بر شاخهٔ بلوطی کهن دیده بود. گمان می‌کرد که
آشیانه‌اش گمان می‌رود در \VUgras{شاخ‌وبرگ} \textsuperscript{(58)}
پنهان باشد و می‌خواست آن را بگیرد.

%%K t09-c2-10-2 p
چالاک چون \VUgras{سنجاب} \textsuperscript{(61)}، از \VUgras{شاخه}
\textsuperscript{(55)} به شاخه بالا رفت، اما چیزی نیافت و تنها توانست
\VUgras{کلاغ} \textsuperscript{(60)} پیری را به پرواز درآورد که بر
\VUgras{شاخهٔ نازک} \textsuperscript{(57)} نشسته بود.

\VUsaut{4.0mm}
\VUfilet{22mm}
